\chapter{Literature Review}
\label{chap:lr}
\chaptermark{Second Chapter Heading}

This chapter provides an overview of existing research related to our topic. The chapter is divided into four main sections, each of which will examine works pertaining to the following
\begin{enumerate}[label=\textbf{Section \arabic*:}, leftmargin=*]
  \item The concept of orthogonal persistence and its implementation options;
  \item PhantomOS and its implementation of the concept of orthogonal persistence;
  \item Details and features of the PhantomOS port to Genode;
  \item Methods of protecting data at rest and their characteristics;
\end{enumerate}

\section{Persistence}

\subsection{Persistence definitions}

Data persistence refers to the period of time during which data exists and is used \cite{PLA}. This is precisely how some of the earliest researchers in the field—Morrison R. and Atkinson M. P.—define this concept. In their work, they identify categories of persistence, divided into two groups: 
\begin{enumerate}[label=\arabic*.]
  \item Provided by the programming language;
  \item Provided by the environment;
\end{enumerate}
Researchers are striving to develop systems in which data usage is independent of the group or category to which the data belongs. 


\subsection{Orthogonal persistence principles}

The concept of orthogonal persistence is based on principles formulated by researchers:
\begin{enumerate}
    \item Independence;
    \item Orthogonality of data types;
    \item Persistence Identification;
\end{enumerate}
In this case, independence means that data persistence does not depend on how the data is manipulated (the user cannot move data between long-term and short-term storage). The orthogonality of data types ensures that there are no cases in which data of any type cannot be persisted. The third principle means that persistent object identification does not relate to the type of such an object. A system that follows all three principles is orthogonally persistent.

\subsection{Advantages of a persistent system}

The desire to create such a system stems from the fact that it offers a number of advantages over conventional systems. These include \cite{Revisited}:
\begin{enumerate}
    \item Increased programming productivity through simplified semantics;
    \item Avoiding ad hoc solutions for data transformation and long-term data storage;
    \item Providing security mechanisms for the entire environment;
    \item Support for gradual evolution;
    \item Automatically preserving referential integrity across the entire computational environment for the entire lifetime of an application.
\end{enumerate}
All of these are, to some extent, the result of simplifying the memory management model for application programmers. This eliminates a huge amount of code that would otherwise be required to handle I/O and data storage. This also includes serialization/deserialization and initialization/deinitialization.

\subsection{The cost of persistence}

All of the advantages outlined above come at a cost. In addition to the general challenges associated with creating and implementing a new system—especially one built on principles that are entirely different from those we are accustomed to—there are also challenges specific to orthogonally persistent systems. These include:
\begin{enumerate}[label=\arabic*.]
    \item The need for a stable object store: This depends largely on the specific implementation method, so we will return to this topic in Section 2.2.
    \item Reduced performance of some applications: Because the programmer does not have full control over the location of the data, access speed may be reduced. A similar problem exists when using virtual memory.
    \item The cost of providing language-independent binding mechanisms.
\end{enumerate}
As will be discussed in Section~2.2, Phantom has solutions to minimize each of these problems. For now, note that point 1 depends largely on the implementation, and it is difficult to generalize information regarding it. Problems similar to point 2 also affect (albeit often to a lesser extent) systems with virtual memory, since the programmer does not have full control over the layout of the data, which can slow down access to it.

\section{PhantomOS}

\subsection{PhantomOS Overview}

PhantomOS is an open-source operating system. The concept was conceived by Dmitry Zavalishin, who also made the primary contribution to the system’s development. Like other modern operating systems, Phantom features multitasking \cite{Habr_MT, Habr_Preempt, Habr_Sched}, a windowing subsystem \cite{Habr_UI}, a network stack, and other familiar attributes. However, for our work, another feature is of greatest interest. PhantomOS provides an orthogonally persistent environment for applications within the Phantom Virtual Machine (PVM), where code is executed in the Phantom programming language. However, it is worth noting that there is also a POSIX compatibility subsystem. Access via arbitrary addresses is prohibited. Phantom already implements subsystems such as:
\begin{itemize}
    \item Kernel itself: threads, synchronization, persistent memory management;
    \item Bytecode virtual machine - running native applications;
    \item POSIX layer - runs Linux-compatible (but not yet persistent) code;
    \item Graphics subsystem - Windows, controls, UI;
    \item Networking (TCP/IP);
    \item Phantom language compiler - the most native userland language;
    \item Java to Phantom translator - work in progress;
    \item Python to Phantom translator - just started;
\end{itemize}
All of this is specified in the documentation \cite{Phantom_docs}.

\subsection{Mechanisms for ensuring persistence}

Dmitry Zavalishin believes that it is impossible to preserve the state of the entire machine, but this is not necessary to achieve persistence \cite{Habr_Persistent_Mem}. Therefore, we will now discuss the mechanisms for ensuring a persistent environment. Persistence in PhantomOS is achieved by page-by-page mapping of the entire memory (PVM) to disk and periodically saving it. The kernel manages the snapshot process. It takes snapshots after first bringing the programs to a state where they are fully represented in memory. That is, all information from the processor registers is also transferred for this purpose. Application programs are blocked only during this brief preparation phase. The disk is switched to read-only mode during this process, but this does not cause a lock throughout the entire process of writing the snapshot to disk thanks to the use of CoW: if a page has been written to the snapshot, nothing prevents access to it; otherwise, a copy is created, which the program operates on, and the original page is moved closer to the front of the save queue \cite{Habr_Persistent_Mem, PM_wiki}. Writing the entire memory to disk with every snapshot is a very costly operation that is unnecessary. To optimize this process, only the increment is written to disk, and preemptive writing is also used. The latter, by the way, serves two purposes:
\begin{itemize}
    \item To speed up the snapshot writing process;
    \item To meet the increased demand for memory during the snapshot;
\end{itemize}

\subsection{System and Persistent Environment features}

One of the goals behind the system’s design is to make life easier for programmers. They no longer need to manage memory; instead, the garbage collector handles this task. Other key features of the system include a global address space and the requirement that objects can only be accessed via a reference. Access via arbitrary addresses is prevented, which has a positive effect on security. It is in this environment that the garbage collector operates. The large amount of virtual memory and the inability to use stop-world algorithms led to a strategy of using two garbage collectors \cite{Habr_GC1}:
\begin{enumerate}
    \item Partial;
    \item Full;
\end{enumerate}
The first is fast, inexpensive, and possibly incomplete. It is assumed that it cleans up garbage in RAM without waiting for objects to be flushed to disk. It runs continuously and is currently implemented based on counting the number of references to an object. The second should be full, but it can run periodically thanks to the presence of the first. At the same time, a stop-world implementation is also possible due to the persistence of garbage: that is, what was garbage in the old snapshot will remain garbage in the new one. Based on this idea, it is proposed to perform garbage collection on the old snapshot, although not the entire snapshot may be used for this, but only a reflection of its memory state. However, there are still a number of problems that need to be solved before its implementation \cite{Habr_GC2}.

\section{OS Phantom to Genode port}

\subsection{Reason for porting}

As previously mentioned, PhantomOS is currently being ported to the Genode framework. This is being done to improve \cite{Antonov_Anton}:
\begin{itemize}
    \item Reliability;
    \item Security;
    \item Interoperability;
\end{itemize}
In this way, we plan to achieve the required quality. One of the main reliability issues—the new kernel with built-in drivers—is planned to be resolved by using kernels supported by Genode, including many microkernels known for their high reliability. It is also proposed to replace new, insufficiently tested components with already proven ones offering the same functionality. This will also resolve the issue of interoperability, as there will be no need to rewrite existing components, and the focus can shift to the system’s logic.  The Genode developers have taken a thorough approach to security. The framework uses a capability-based security system. Microkernels are often used as the system kernel. This makes it possible to isolate components to a significant degree and configure access to resources more flexibly, allowing for a minimal TCB. This results in a minimal attack surface, which has an extremely positive effect on reliability.

\subsection{Affected components}

Initially, work was carried out to port the system \cite{Antonov_Anton}. Essentially, the task involves migrating the persistent environment to the framework. In addition to PVM, work was done on HAL, libc, and the physical memory allocation mechanism. Some functions were replaced with placeholders. The result of this work was a working isomem project \cite{GitHub_isomem}. It also includes a video placeholder in the form of a window subsystem with limited functionality. It is configured and runs as a separate, standalone Genode component. Subsequently, work was carried out on developing a persistent network subsystem for the \cite{Brisilin_Anton} port and implementing the WASM bytecode runtime \cite{Samburskiy_Kirill}. The snapshot creation mechanism was also modified, resulting in the creation of the snapper project \cite{GitHub_Snapper}.

\subsection{Snapper}

The current version of Snapper (Snapper 2.0) is designed to address the issues encountered in the implementation of the previous mechanism—the “superblock” \cite{FOSDEM_Snapper, FOSDEM_Snapper_paper}. It improves disk space efficiency when storing multiple snapshots. Previously, snapshots were stored as monolithic objects in memory, which resulted in a large number of data copies when storing multiple snapshots. For unchanged pages, Snapper adds a link in the new generation without copying all the data again. The ability to recover from failures is ensured by using file systems for which such tools already exist, for example: ext4 and fsck. Using a file system to store snapshots can be more convenient during development and debugging. Integrity checks have been added.  In addition, the system has become configurable in terms of reliability through the introduction of a parameter called redundancy. If the specified number of links to a file is exceeded, a copy of the file is created, and all generations using this file are linked to both the original file and its copy. This allows an experienced administrator to independently configure the balance between reliability and storage usage.

\section{Data protection at rest}

Snapper has helped improve snapshot security by adding integrity checks, but neither it nor any other component ensures secure storage. They are stored on disk in plaintext, which can pose a serious security risk in many device usage scenarios.

\subsection{Threats to stored data}

Data at rest refers to any data stored on persistent storage media that is not currently being actively transmitted or processed. It is static and vulnerable to threats related to physical or logical access. Sensitive data may be stored on the disk and could be read or overwritten. In the context of PhantomOS, the situation is complicated by the fact that a snapshot is taken regardless of what data is in memory, and secret data may end up in it, even if the application level of the program provides for its protection using its own means \cite{SNIA_PM_Security}. This is similar to the problem of protecting virtual machine snapshots \cite{VM_Security}.

\subsection{Ways to protect stored data}

Data protection methods must align with the threat model, which, in turn, depends directly on the device’s usage scenarios. This subsection provides only a brief overview of such methods, without a detailed explanation of why a particular method might be chosen.
In this case, we are discussing storage encryption. We do not consider options involving robust physical and/or logical access controls. Encryption can be classified by scope \cite{NIST_SP800111}:
\begin{itemize}
    \item Full disk encryption;
    \item Volume and virtual disk encryption;
    \item File and folder encryption;
\end{itemize}
and by level \cite{CryptoFS}:
\begin{itemize}
    \item Block level;
    \item File system level (generalized, divided into several points regarding the number and removal of file systems);
    \item Application level;
\end{itemize}
These are the levels used at the OS level and above. Each higher level leaves the metadata of the lower level exposed. In addition to these, implementation at a low-level hardware level is possible, which would be transparent to the OS. Due to the absence in Genode of verified systems of this type or tools for their creation, this level will not be considered.

\subsection{Current work}

In this project, we will focus on implementing an environment for the secure storage of snapshots. This will be a block-level encrypted storage solution, which will enhance security compared to file-system-level encryption and higher, by hiding metadata and making it more difficult to extract system information. This is also facilitated by the fact that Genode already has a library \cite{GitHub_tresor} for this purpose and a test component \cite{GitHub_file_vault, Page_file_vault}. 

Tresor performs AES encryption on 4K-sized blocks. The virtual block device key is encrypted with a master key, which is managed by the Trust Anchor. Copy-on-write (CoW) is used during writing, ensuring the atomicity of the operation and improving the ability to recover from failures. The component uses SHA-256 for data integrity checking.
